\chapter{Классификация калибровочно-ковариантного старо-нового коннектора}
\label{ch:v8-baryon-c0-old-new-gauge-covariant-connector-classification-gate}

Предыдущий гейт показал, что внутренний паулиевский блок не уничтожает
центральный остаток $\mathbb C^2$. Ищем отображение
$V:H_{\mathrm n}\to H_{21}$, которое не обязано быть gauge-инвариантным
по отдельности, но должно порождать замкнутый ковариантный кадр скачков.

Пусть $G_1,\ldots,G_{11}$ --- генераторы $SU(3)\times SU(2)$, а $Y$ ---
гиперзаряд. Минимальный заряженный класс задаётся системой
\begin{equation}
 G_jV=0,
 \qquad (Y+I_{21})V=0,
 \qquad \Gamma_{21}V+V\Gamma_{\mathrm n}=0.
 \label{eq:v8-baryon-c0-old-new-covariant-map-system}
\end{equation}
Её точная размерность и базис равны
\begin{equation}
 \dim_{\mathbb C}\mathcal V_{-1}^{\rm odd}=3,
 \qquad
 \mathcal V_{-1}^{\rm odd}
 =\operatorname{span}_{\mathbb C}
 \{|e_R^{(s)}\rangle\langle s_0|,
 |e_R^{(t,1)}\rangle\langle a_0|,
 |e_R^{(t,2)}\rangle\langle a_0|\}.
 \label{eq:v8-baryon-c0-old-new-covariant-map-basis}
\end{equation}
Иными словами,
\begin{equation}
 V(z)=z_0|e_R^{(s)}\rangle\langle s_0|
 +z_1|e_R^{(t,1)}\rangle\langle a_0|
 +z_2|e_R^{(t,2)}\rangle\langle a_0|.
 \label{eq:v8-baryon-c0-old-new-general-connector}
\end{equation}
Стандартная Real-совместимость допускает вещественного представителя с
точностью до общего масштаба, поэтому множество направлений остаётся
неединственным:
\begin{equation}
 [z_0:z_1:z_2]\in\mathbb{RP}^2.
 \label{eq:v8-baryon-c0-old-new-real-projective-family}
\end{equation}

\section{Минимальная ковариантная орбита}

Для любого ненулевого вещественного $V$ положим
\begin{equation}
 C_1=\begin{pmatrix}0&V\\V^*&0\end{pmatrix},
 \qquad
 C_2=\begin{pmatrix}0&iV\\-iV^*&0\end{pmatrix}.
 \label{eq:v8-baryon-c0-old-new-hermitian-quadratures}
\end{equation}
Нетривиальные генераторы коммутируют с обоими направлениями, а гиперзаряд
вращает их пару:
\begin{equation}
 i[Y,C_1]=-C_2,
 \qquad i[Y,C_2]=C_1.
 \label{eq:v8-baryon-c0-old-new-u1-frame-rotation}
\end{equation}
Оба направления нечётны и образуют Real-замкнутый span:
\begin{equation}
 \{\Gamma_{23},C_1\}=\{\Gamma_{23},C_2\}=0,
 \qquad J_0C_1J_0^{-1}=C_1,
 \qquad J_0C_2J_0^{-1}=-C_2.
 \label{eq:v8-baryon-c0-old-new-grading-real-closure}
\end{equation}
Для нормированного rank-one представителя новая пара следово ортогональна
всему блочно-диагональному 45-кадру и
\begin{equation}
 \mathcal F_{47}=\mathcal F_{45}\oplus
 \operatorname{span}_{\mathbb R}\{C_1,C_2\},
 \qquad
 K_{47}=K_{45}\oplus2I_2,
 \qquad \rank K_{47}=47.
 \label{eq:v8-baryon-c0-old-new-frame-47-metric}
\end{equation}

\section{Уничтожение второго центрального направления}

В ортонормированном центральном базисе
$P_{21}/\sqrt{21},P_{\mathrm n}/\sqrt2$ вклад двух новых скачков равен
\begin{equation}
 D_{\rm cen}=
 \begin{pmatrix}
 2/21&-\sqrt{42}/21\\
 -\sqrt{42}/21&1
 \end{pmatrix},
 \qquad \rank D_{\rm cen}=1,
 \qquad \ker D_{\rm cen}=\mathbb C I_{23}.
 \label{eq:v8-baryon-c0-old-new-central-dirichlet-matrix}
\end{equation}
Следовательно, для любого положительного коэффициента $\gamma$,
\begin{equation}
 \mathcal L_{47}=\mathcal L_{45}
 +\gamma\sum_{r=1}^{2}
 \left(C_r(\cdot)C_r-\frac12\{C_r^2,\cdot\}\right)
 \label{eq:v8-baryon-c0-old-new-connected-generator}
\end{equation}
имеет
\begin{equation}
 \operatorname{Fix}(\mathcal L_{47})=\mathbb C I_{23}.
 \label{eq:v8-baryon-c0-old-new-scalar-fixed-algebra}
\end{equation}

\begin{theorem}[Минимальная gauge-ковариантная динамическая сшивка]
Двух Hermitian quadratures одного ненулевого заряженного нечётного
коннектора достаточно, чтобы сократить неподвижную алгебру с
$\mathbb C^2$ до $\mathbb C I_{23}$. Двухмерность вещественной орбиты
минимальна из-за ненулевого $U(1)$-заряда.
\end{theorem}

\begin{nogocriterion}[Ковариантность не выбирает ветвь и скорость]
Допустимые Real-направления образуют $\mathbb{RP}^2$, а положительный
коэффициент $\gamma$ остаётся свободным. Поэтому существование примитивного
47-скачкового продолжения не выводит конкретный коннектор, его скорость или
физическое значение $c_0$ из текущего parent action.
\end{nogocriterion}

\begin{protocol}\begin{enumerate}
\item комплексная размерность допустимых нечётных карт: \textbf{$3$};
\item Real-проективное семейство: \textbf{$\mathbb{RP}^2$};
\item минимальный вещественный gauge-кадр: \textbf{$2$ направления};
\item расширенный frame rank: \textbf{$47$};
\item центральная дирахлеевская матрица: \textbf{ранг $1$};
\item полная неподвижная алгебра: \textbf{$\mathbb C I_{23}$};
\item примитивность при $\gamma>0$: \textbf{да};
\item выбор ветви и $\gamma$: \textbf{не выведен};
\item следующий гейт: \textbf{parent-селектор multiplicity-направления и
скорости коннектора}.
\end{enumerate}\end{protocol}

Машинный сертификат сохранён в
\path{s2t_v8_baryon_c0_old_new_gauge_covariant_connector_classification_gate_results.json}.